% Noether P09 Korean producer draft — UNCHECKED
% T07-U59; ED0002 lines7051--7073; source 1729 LF B / 6FD68A32E861265512D31C1CA834F4117D396364F9D59BC85B257F25CF48E472
% Pointer v009 / ED0002: B06BE3530D9CF2E82B56FDBA7FE41D5D044DF2425DFA2A059D4939EAA2F7A6C2 / C9A125167ACB33D914EE4374B65AE7CDF0052F568371B8B77B720EA178ABF0E3
% Translation only; induction, side-condition, and theorem wording remain checker-held.

$\mathfrak G$의 나머지 양 $\vartheta$들은 계 $\Theta$를 이루며, 이 계가 $\mathfrak G$의 모든 양의 유리기저임을 보이겠다. 실제로 가정에 따라 각 양 $\vartheta$는 $\Theta$의 앞선 양들에 유리적으로 종속되지 않는다. 또한 귀납 논법\srcfn{***)}{Zermelo, 앞의 논문 § 1 참조.}은 모든 관계 $z=\varphi(\xi)$가 관계 $z=\psi(\vartheta)$를 수반함을 보여 준다. 그렇지 않다면 $\vartheta$들로 유리적으로 나타낼 수 없는 첫 원소 $z_0=\varphi_0(\xi_1,\ldots,\xi_t)$가 존재해야 하지만, 그보다 앞선 $\xi_1,\ldots,\xi_t$는 $\vartheta$들의 유리함수이므로 모순이 생긴다. $\Theta$는 II에 따라 동시에 모든 수의 유리기저이다. III와 I에 따라 $\mathfrak G$는 $\Theta$와 더불어 $\vartheta$의 모든 정수적 다항식으로 이루어진 정역 $[\Theta]$를 부분정역으로 포함하며, IV에 따라 $[\Theta]$는 어떠한 분수형 대수수도 포함하지 않는다. 따라서 모든 수의 유리기저 $\Theta$는, 그로부터 생기는 정역 $[\Theta]$가 어떠한 분수형 대수수도 포함하지 않는다는 부가조건을 만족한다.\srcfn{*)}{이것이 실제 조건이라는 사실은, 예를 들어 두 양 $\eta,\frac1{2\sqrt{\eta}}$는 기저 원소로 허용되지 않지만 $\eta,\frac1{\sqrt{\eta}}$는 허용된다는 데서 드러난다.} 즉 $\Theta$는 모든 수에 대한 \emph{부가조건을 만족하는 유리기저}이다. 그러므로 § 2와 유사하게 다음이 성립한다. \emph{각 영역 $\mathfrak G$는 부가조건을 만족하는 유리기저를 이용하여 구성할 수 있으며, 이때 $\mathfrak G$는 정역 $[\Theta]$를 부분정역으로 포함한다.}
